A live equity portfolio poses three problems: what to hold, when to trade, and whether the strategy holds up beyond
the observed market. We built one system for all three and backtested it on 373 US equities, calibrated on 2014--2024
and traded January 2025 to September 2026. Holdings maximized a utility of shares held under limits on budget,
single-stock weight, portfolio volatility and turnover, with each stock's preference weight a bounded score of its
trend, market loading and the market's recent direction. We compared seven weighting rules (four consumer-theory
utilities, three standard portfolios) across rules picking stocks by score, by bandit or at random. A model-predictive
controller forecast portfolio value with a multi-asset geometric Brownian motion and re-planned at its horizon, when
the value left its forecast band, or at a drawdown stop. We also ran the system on simulated price paths (synthetic
futures) from per-asset hidden Markov generators with jumps, linked by market and sector factors. On real prices,
terminal wealth spread more across selection rules than weighting rules (median spreads 0.52 and 0.14 of initial
wealth). Band exits and drawdown stops triggered a fifth to three fifths of re-plans, which, against re-planning every
21 trading days, cut the worst drawdown in most rule combinations without consistently raising wealth. In the
synthetic futures with the deepest market declines, every rule combination, random picks included, ended above
simulated SPY and the equal-weight universe in the median, so the margin needed no stock selection. A market hedge (a
market-loading cap that tightened as the market weakened, plus zero- or negative-loading stocks) cost wealth on real
prices but added to it in the deepest synthetic declines without jump bursts. The synthetic market's persistent trends
favored score-based over random selection more than the real window did.
